\documentclass[a4paper,%
11pt,%
DIV12,
headsepline,%
headings=normal,
]{scrartcl}

\usepackage{bopc}
\usepackage{csvsimple}
\usepackage{booktabs}
\usepackage{float}
\usepackage{graphicx}
\usepackage{caption}

%% ==================== ADAPT ACCORDINGLY ==================== %%

% args: first name, last name, matriculation number
\setAuthorOne{Jane}{Doe}{12012345}
\setAuthorTwo{John}{Roe}{12112345}
\setAuthorThree{Max}{Min}{14312345}

% args: group size, group number (TUWEL)
\setGroup{3}{X}

%% =========================================================== %%

%% Convenience path macros
\newcommand{\resultpath}{../results}
\newcommand{\plotpath}{../plots}

\begin{document}

\maketitlepage

%% =================== EXERCISE 6 ONLY ==================== %%

\section{Exercise 6}

%% ------------------------------------------------------------------
\subsection{Convert OpenMP code to CUDA}

% No benchmarking content for subtask 1.
% Add code-level description / explanation here if required.

%% ------------------------------------------------------------------
\subsection{Running time analysis}
\label{sec:ex6-subtask2}

%
% Measurement setup:
%   GPU: juliaset_gpu, default block size 16x16, nrep=5
%   Resolutions: 1000, 2000, 4000, 8000, 16000 (square images)
%   Hardware: Tesla T4, CUDA 12.4
%   Rep 0 includes CUDA context initialisation; reps 1-4 are steady-state.
%

\begin{figure}[H]
  \centering
  \includegraphics[width=0.85\textwidth]{\plotpath/subtask2.png}
  \caption{GPU kernel time vs.\ resolution (default $16{\times}16$ block size,
    $\texttt{nrep}=5$).
    \emph{Red}: maximum time across 5 reps (rep~0, includes CUDA context
    initialisation).
    \emph{Blue}: average steady-state time (reps~1--4).}
  \label{fig:subtask2}
\end{figure}

\begin{table}[H]
  \centering
  \caption{Subtask~2 summary: GPU timing vs.\ resolution.
           \texttt{avg\_steady\_s} excludes rep~0 (init-dominated).}
  \label{tab:subtask2}
  \begin{tabular}{rrrr}
    \toprule
    Resolution & Avg steady (s) & Max (s) & Speedup vs.\ 1\,000 \\
    \midrule
    \csvreader[
      head to column names,
      late after line=\\,
    ]{\resultpath/subtask2/subtask2_summary.csv}{}%
    {\resolution & \avgsteadys & \maxs & \speedupvsk}%
    \bottomrule
  \end{tabular}
\end{table}

% TODO: discussion
% KEYWORDS:
%   - steady-state time scales ~4x when resolution doubles (pixels 4x):
%     confirms per-pixel cost dominates at large N
%   - max (rep 0) is nearly flat across resolutions (~250-760 ms):
%     dominated by fixed CUDA context init + cudaMalloc, not pixel count
%   - asymptote: as N -> inf, avg/max -> 1 (init amortised away)
%   - why max > avg: warm-up / context init on first kernel launch per process

%% ------------------------------------------------------------------
\subsection{Impact of block size}
\label{sec:ex6-subtask3}

%
% Measurement setup:
%   GPU: juliaset_gpu, r=20000x20000, nrep=5
%   Block sizes: (1,1), (32,1), (1,32), (128,1), (1024,1)
%   Profiler: ncu (Nsight Compute 2024.1.1) — nvprof unsupported on CC 7.5
%   Metrics: achieved occupancy, global store throughput (L1 pipe),
%            warp thread instruction efficiency
%

\begin{figure}[H]
  \centering
  \includegraphics[width=0.85\textwidth]{\plotpath/subtask3.png}
  \caption{Average steady-state GPU time per block size
    ($r=20\,000{\times}20\,000$, $\texttt{nrep}=5$, reps~1--4).
    Best configuration highlighted in green.}
  \label{fig:subtask3}
\end{figure}

\begin{table}[H]
  \centering
  \caption{Subtask~3 summary: timing and ncu profiler metrics per block size.
           GST = global store throughput (L1 texture pipe).
           Warp eff = active threads per warp-instruction (max 32).}
  \label{tab:subtask3}
  \begin{tabular}{lrrrrr}
    \toprule
    Block & Threads & Avg (s) & Occupancy (\%) & GST (GB/s) & Warp eff \\
    \midrule
    \csvreader[
      head to column names,
      late after line=\\,
    ]{\resultpath/subtask3/subtask3_summary.csv}{}%
    {\block & \threads & \avgsteadys & \occupancypct & \gstGBs & \warpeff}%
    \bottomrule
  \end{tabular}
\end{table}

% TODO: discussion
% KEYWORDS:
%   - indexing julia_set[x*res_y + y]: consecutive y = consecutive addresses
%     => (1,32) threads vary y within warp => coalesced global store (1 transaction/warp)
%     => (32,1) threads vary x => stride res_y=20000 => 32 transactions/warp (non-coalesced)
%   - (1,32) wins despite lower L1-pipe throughput metric: high GST for (32,1)
%     reflects wasted cache-line fill traffic, not useful DRAM bandwidth
%   - warp size 32: (1,1) uses 1/32 warp lanes => warp_eff=1 (3% utilisation);
%     all other configs fill full warps
%   - (1,1) is 7x slower: both 1/32 lane utilisation AND non-coalesced stores
%   - (128,1) has highest occupancy (81%) but is 21% slower than (1,32):
%     occupancy alone does not predict performance when memory coalescing limits
%   - global loads = 0: kernel only WRITES julia_set; c and params in registers
%   - achieved vs theoretical occupancy: (1,32) 49% actual vs 100% theoretical
%     (limited by warp scheduling, not register/smem)

%% ------------------------------------------------------------------
\subsection{Running time: CPU vs.\ GPU}
\label{sec:ex6-subtask4}

%
% Measurement setup:
%   GPU: juliaset_gpu, block (1,32), nrep=5
%   CPU: juliaset_cpu, OMP_NUM_THREADS=<nproc on tesla node>
%   Resolutions: 1000, 2000, 4000, 8000, 16000
%   nproc recorded at job runtime: see results/subtask4/job.out
%

\begin{figure}[H]
  \centering
  % subtask4.png generated once results/subtask4/ is populated
  \includegraphics[width=0.85\textwidth]{\plotpath/subtask4.png}
  \caption{CPU vs.\ GPU average steady-state time vs.\ resolution
    (log--log scale, $\texttt{nrep}=5$, reps~1--4).
    GPU uses best block $(1{\times}32)$; CPU uses \texttt{OMP\_NUM\_THREADS=nproc}.
    Dashed vertical line marks the crossover resolution (if within range).}
  \label{fig:subtask4}
\end{figure}

\begin{table}[H]
  \centering
  \caption{Subtask~4 summary: CPU vs.\ GPU average steady-state time.
           Speedup $= t_\text{CPU} / t_\text{GPU}$.}
  \label{tab:subtask4}
  \begin{tabular}{rrrr}
    \toprule
    Resolution & GPU avg (s) & CPU avg (s) & Speedup \\
    \midrule
    \csvreader[
      head to column names,
      late after line=\\,
    ]{\resultpath/subtask4/subtask4_summary.csv}{}%
    {\resolution & \gpuavgs & \cpuavgs & \speedup}%
    \bottomrule
  \end{tabular}
\end{table}

% TODO: discussion
% KEYWORDS:
%   - crossover point: GPU overhead (cudaMalloc + memcpy ~250 ms fixed cost)
%     dominates at small N; CPU wins below crossover
%   - at large N: GPU scales ~O(N^2) with high parallelism => GPU wins
%   - CPU baseline depends on OMP_NUM_THREADS (record exact nproc value);
%     more threads => lower CPU time => pushes crossover to larger N
%   - GPU has constant overhead regardless of resolution:
%     cudaMalloc + H2D memcpy ~0 here (output only) + D2H memcpy ~O(N^2)
%   - both axes log-scale: straight-line slopes reveal power-law scaling
%   - if no crossover in [1000,16000]: note GPU is always faster or always slower
%     and extrapolate where crossing would occur

\end{document}
